\capitulo{2}{Objetivos del proyecto}

El objetivo principal del proyecto es diseñar e implementar un sistema de Retrieval-Augmented Generation (RAG) que permita enriquecer las respuestas de un modelo de lenguaje mediante la recuperación y utilización de información contextual proveniente de una colección documental. 

\section{Objetivos específicos}
Los objetivos identificados en el presente trabajo fin de grado son los siguientes:
\begin{itemize}
\tightlist
\item
Preprocesar la colección documental.
	\begin{itemize}
		\item Segmentar los documentos en fragmentos.
		\item Normalizar y limpiar el texto.
		\item Generar \textit{embeddings} para representar semánticamente los fragmentos.
	\end{itemize}
	\item
Implementar y gestionar un índice vectorial.
	\begin{itemize}
		\item Almacenar los \textit{embeddings} en una base de datos vectorial.
		\item Optimizar consultas de similitud para recuperar contexto relevante.
	\end{itemize}
\item Desarrollar la capa de interacción con el usuario
	\begin{itemize}
		\item Transformar la consulta en su representación vectorial.
		\item Recuperar los fragmentos más relevantes del índice.
		\item Integrar los fragmentos en el prompt para enriquecer la respuesta del LLM.
	\end{itemize}
\end{itemize}

\section{Objetivos personales}

\begin{itemize}
\tightlist
\item Adquirir conocimientos sobre los sistemas RAG, comprendiendo sus fundamentos, componentes y aplicaciones en el ámbito de la inteligencia artificial.
\item Profundizar en el estudio de los modelos de lenguaje de gran tamaño (LLM).
\item Desarrollar habilidades prácticas en la implementación de sistemas basados en RAG y LLM.
\item Aprender a aplicar los sistemas RAG y los LLM en la resolución de problemas reales.
\item Aplicar los conocimientos adquiridos durante la carrera.
\end{itemize}
